\begin{examplebox}
    \textbf{\textcolor{CExample}{例一（基础：四次首尾对称方程降为二次）}}
    \textbf{\textcolor{CExample}{题目：}}
    解方程 $t^4-3t^3+4t^2-3t+1=0$。
    \begin{description}[style=nextline, leftmargin=2.8em, labelsep=0.5em, itemsep=0.5em]
        \item[\textbf{第一步（识别并确认 $t\ne0$）：}]
            系数为 $1,-3,4,-3,1$，首尾成对相同；且 $t=0$ 不是根。
        \item[\textbf{第二步（除以 $t^2$）：}]
            \[
               \left(t^2+\frac1{t^2}\right)-3\left(t+\frac1t\right)+4=0.
            \]
        \item[\textbf{第三步（整体换元，验范围）：}]
            令 $u=t+1/t$，则 $|u|\ge2$，原式变为
            \[
              u^2-2-3u+4=(u-1)(u-2)=0.
            \]
            所以 $u=1$ 或 $u=2$；但 $u=1$ 不在允许范围内。
        \item[\textbf{第四步（回代）：}]
            \[
              t+\frac1t=2\Longleftrightarrow t^2-2t+1=0
              \Longleftrightarrow\boxed{t=1}.
            \]
            \textbf{理由：} 换元解 $u$ 不等于原方程解完，必须回代检查。
    \end{description}

    \end{examplebox}
\clearpage
\begin{examplebox}
    \textbf{\textcolor{CExample}{例二（最值：选择更短的二次倒数和换元）}}
    \textbf{\textcolor{CExample}{题目：}}
    设实数 $t\ne0$，求
    \[
      P=t^4+\frac1{t^4}-3\left(t^2+\frac1{t^2}\right)
    \]
    的最小值。
    \begin{description}[style=nextline, leftmargin=2.8em, labelsep=0.5em, itemsep=0.52em]
        \item[\textbf{第一步（只含偶数次，直接换二次整体）：}]
            令 $z=t^2+1/t^2$。由 $(t-1/t)^2\ge0$，得 $z\ge2$。
        \item[\textbf{第二步（降次并整理）：}]
            \[
              t^4+\frac1{t^4}=z^2-2,\qquad
              P=z^2-3z-2.
            \]
        \item[\textbf{第三步（在允许范围内比较）：}]
            \begin{center}
\begin{tikzpicture}[font=\small]
\begin{axis}[width=11.8cm,height=6cm,axis lines=middle,
 xmin=0,xmax=4.1,ymin=-5.1,ymax=3.5,xtick={1.5,2,3},ytick={-4.25,-4},
 yticklabels={},xlabel={$z$},ylabel={$P$},clip=false]
\addplot[gray,dashed,domain=.15:2,samples=60] {x^2-3*x-2};
\addplot[CStatement,very thick,domain=2:4,samples=80] {x^2-3*x-2};
\addplot[only marks,mark=o,gray,mark size=2.5pt] coordinates {(1.5,-4.25)};
\addplot[only marks,CExample,mark size=3pt] coordinates {(2,-4)};
\node[anchor=east,text=gray] at (axis cs:1.25,-2.4) {顶点不允许};
\node[anchor=west,text=CExample] at (axis cs:2.2,-4.35) {$(2,-4)$ 可取};
\node at (axis cs:1.6,2.6) {$P=z^2-3z-2$};
\node at (axis cs:1.6,1.5) {实线只保留 $z\ge2$};
\end{axis}
\end{tikzpicture}
\end{center}

            \[
              P+4=z^2-3z+2=(z-1)(z-2)\ge0,
            \]
            因为 $z\ge2$，两因式均不小于零。
        \item[\textbf{第四步（验取等条件）：}]
            $P=-4$ 当且仅当 $z=2$，即 $(t-1/t)^2=0$，所以 $t=\pm1$。故
            \[
               \boxed{P_{\min}=-4\quad (t=\pm1)}.
            \]
            \textbf{理由：} 直接设 $z=t^2+t^{-2}$ 比先设 $u=t+t^{-1}$ 少一次替换。
    \end{description}

    \end{examplebox}
\clearpage
\begin{examplebox}
    \textbf{\textcolor{CExample}{例三（圆锥曲线迁移：抛物线两条垂线的中点距离）}}
    \textbf{\textcolor{CExample}{题目：}}
    对抛物线 $y^2=8x$，焦点为 $F(2,0)$，准线为 $l:x=-2$。过焦点的两条互相垂直的直线与抛物线分别交于 $A,B$ 和 $D,E$，两条焦点弦中点为 $M,N$。已知通过联立方程与根的平均值，设非零斜率为 $t$ 时有
    \[
      M\left(2+\frac4{t^2},\frac4t\right),\qquad
      N(2+4t^2,-4t).
    \]
    求 $|MN|$ 的最小值，并判断取最小时 $MN$ 与准线的关系。
    \begin{center}
\begin{tikzpicture}[font=\small]
\begin{axis}[width=10cm,height=11.3cm,scale only axis,axis equal image,axis lines=middle,
 xmin=-3,xmax=13,ymin=-11,ymax=11,xtick=\empty,ytick={-4,4},
 xlabel={$x$},ylabel={$y$},clip=false]
\addplot[CStatement,very thick,domain=-10.2:10.2,samples=180,variable=\s] ({\s^2/8},{\s});
\addplot[CTrap,dashed,thick] coordinates {(-2,-10.5) (-2,10.5)};
\addplot[gray,domain=.34314575:11.65685425] {x-2};
\addplot[gray,domain=.34314575:11.65685425] {-x+2};
\addplot[CExample,very thick] coordinates {(6,-4) (6,4)};
\addplot[only marks,CExample,mark size=2.5pt] coordinates {(6,4) (6,-4)};
\addplot[only marks,CStatement] coordinates {(2,0) (11.65685425,9.65685425) (.34314575,-1.65685425) (11.65685425,-9.65685425) (.34314575,1.65685425)};
\draw (axis cs:2.4,.4)--(axis cs:2.8,0)--(axis cs:2.4,-.4);
\node[anchor=west] at (axis cs:6.35,4.35) {$M(6,4)$};
\node[anchor=west] at (axis cs:6.35,-4.35) {$N(6,-4)$};
\node[anchor=west] at (axis cs:3.3,-.65) {$F$};
\node[anchor=west] at (axis cs:6.4,1.5) {$|MN|=8$};
\node[anchor=south] at (axis cs:11.65,9.9) {$A$};
\node[anchor=north] at (axis cs:11.65,-9.9) {$D$};
\node[anchor=east] at (axis cs:0,-2.3) {$B$};
\node[anchor=east] at (axis cs:0,2.3) {$E$};
\node[rotate=90,text=CTrap] at (axis cs:-3.5,6) {$l:x=-2$};
\node at (axis cs:6,10.3) {$y^2=8x$，$t=1$};
\end{axis}
\end{tikzpicture}
\end{center}

    图示 $t=1$ 的等号状态：两条焦点弦互相垂直，$MN$ 与准线同为竖直。一般 $t$ 的最小性仍由下列计算证明。
    \begin{description}[style=nextline, leftmargin=2.8em, labelsep=0.5em, itemsep=0.55em]
        \item[\textbf{第一步（先化为坐标差）：}]
            \[
              \begin{aligned}
              |MN|^2&=\left(4t^2-\frac4{t^2}\right)^2
                         +\left(-4t-\frac4t\right)^2\\
                 &=16\left[\left(t^2-\frac1{t^2}\right)^2
                    +\left(t+\frac1t\right)^2\right].
              \end{aligned}
            \]
        \item[\textbf{第二步（发现倒数和整体）：}]
            令 $u=t+1/t$，则 $u^2\ge4$，且
            \[
              \left(t^2-\frac1{t^2}\right)^2
              =\left(t-\frac1t\right)^2\left(t+\frac1t\right)^2
              =(u^2-4)u^2.
            \]
        \item[\textbf{第三步（用二次式证明最小）：}]
            令 $z=u^2\ge4$，则
            \[
              |MN|^2=16z(z-3),\quad
              |MN|^2-64=16(z-4)(z+1)\ge0.
            \]
            因此 $|MN|\ge8$。
        \item[\textbf{第四步（验等号并检查方向）：}]
            $z=4$ 等价于 $t=1$ 或 $t=-1$。图中 $t=1$ 时两点横坐标相同、纵坐标相差 $8$；$t=-1$ 时两点上下交换，方向和长度不变。
            \[
              \boxed{|MN|_{\min}=8,\qquad MN\parallel l.}
            \]
            \textbf{理由：} 最小值为 $8$，不是 $0$；两点横坐标相等时，纵坐标仍相差 $8$。
    \end{description}
    \textbf{考场更快的比较法：} $|MN|\ge|\Delta y|=4|t+1/t|\ge8$，等号在 $t=\pm1$ 时两步同时成立。整体换元法适合整理复杂平方，坐标差下界法通常更快。

    \end{examplebox}
\clearpage
\begin{examplebox}
    \textbf{\textcolor{CExample}{例四（倒数差：根据表达式选换元）}}
    \textbf{\textcolor{CExample}{题目：}}
    解方程
    \[
       t^2+\frac1{t^2}-3\left(t-\frac1t\right)=0,\qquad t\ne0.
    \]
    \begin{description}[style=nextline, leftmargin=2.8em, labelsep=0.5em, itemsep=0.48em]
        \item[\textbf{第一步（令差为一个整体）：}]
            令 $v=t-1/t$，则 $v\in\mathbb R$，$t^2+1/t^2=v^2+2$。
        \item[\textbf{第二步（解二次方程）：}]
            \[
              v^2-3v+2=(v-1)(v-2)=0\Longrightarrow v=1\text{ 或 }2.
            \]
        \item[\textbf{第三步（回代）：}]
            若 $v=1$，则 $t^2-t-1=0$，得 $t=\dfrac{1\pm\sqrt5}{2}$；若 $v=2$，则 $t^2-2t-1=0$，得 $t=1\pm\sqrt2$。
        \item[\textbf{第四步（列全答案）：}]
            \[
               \boxed{t=\frac{1\pm\sqrt5}{2}\quad\text{或}\quad t=1\pm\sqrt2.}
            \]
            \textbf{理由：} 若误用 $|v|\ge2$，会错误舍去 $v=1$ 的两个根。
    \end{description}

    \end{examplebox}
\clearpage
\begin{examplebox}
    \textbf{\textcolor{CExample}{例五（指数迁移：先把指数看作倒数变量）}}
    \textbf{\textcolor{CExample}{题目：}}
    解 $4^x+4^{-x}=2^x+2^{-x}+4$。
    \begin{description}[style=nextline, leftmargin=2.8em, labelsep=0.5em, itemsep=0.48em]
        \item[\textbf{第一步（设正数）：}]
            设 $t=2^x>0$，则 $4^x=t^2$、$4^{-x}=t^{-2}$。
        \item[\textbf{第二步（再设倒数和）：}]
            令 $u=t+1/t\ge2$，则 $u^2-2=u+4$，即
            \[
               (u-3)(u+2)=0.
            \]
            因 $u\ge2$，取 $u=3$。
        \item[\textbf{第三步（回代到 $x$）：}]
            $t+1/t=3\Longrightarrow t^2-3t+1=0$，故 $t=(3+\sqrt5)/2$ 或 $t=(3-\sqrt5)/2$，均为正数。
            \[
              \boxed{x=\log_2\frac{3+\sqrt5}{2}\quad\text{或}\quad x=\log_2\frac{3-\sqrt5}{2}.}
            \]
            \textbf{理由：} 指数换元使 $t>0$，因而 $u$ 只取 $[2,+\infty)$。
    \end{description}
\end{examplebox}

\clearpage
